\chapter{Giới thiệu}
\label{ch:introduction}

Chương này thiết lập bối cảnh nghiên cứu cho SENTINEL, một kiến trúc nhận thức hai tầng nhằm tự động hóa việc phát hiện và phản hồi sự cố trong các Trung tâm Điều hành An ninh (SOC) hiện đại. Thông qua quá trình phân tích tình trạng quá tải cảnh báo (alert fatigue), sự hạn chế của các hệ thống dựa trên tập luật tất định (deterministic rule blindness), và các rào cản tính toán của Mô hình Ngôn ngữ Lớn (LLMs), chương này trình bày cơ sở thiết kế và xác thực một đường ống (pipeline) phát hiện lai bảo mật. Các nội dung chính bao gồm mục tiêu, câu hỏi nghiên cứu, phạm vi, và phương pháp luận định lượng.

\section{Bối cảnh và Động lực}
Các hệ thống mạng doanh nghiệp hiện đại đang có xu hướng dịch chuyển mạnh mẽ sang kiến trúc điện toán đám mây (cloud-native) với ranh giới an ninh ngày càng trở nên mờ nhạt, dẫn đến sự mở rộng đáng kể của bề mặt tấn công. Các tác nhân đe dọa sử dụng một phổ rộng các kỹ thuật, từ mã độc (malware) thông thường cho đến các chiến dịch tấn công đa giai đoạn tinh vi~\cite{mlids2016}. Bất chấp sự đa dạng này, các lực lượng phòng thủ luôn phải đối mặt với một thách thức chung: khối lượng khổng lồ của dữ liệu giám sát an ninh (telemetry) cần được sàng lọc. Các SOC hiện nay tiến hành tổng hợp nhật ký sự kiện (log) từ Tường lửa Thế hệ mới (NGFW), Hệ thống Phát hiện/Ngăn ngừa Xâm nhập (IDS/IPS), và Hệ thống Phát hiện và Phản hồi Điểm cuối (EDR) vào các nền tảng Quản lý Thông tin và Sự kiện An ninh (SIEM) tập trung nhằm mục đích chuẩn hóa và tương quan dữ liệu~\cite{siem2021}. Tuy nhiên, các cơ chế lọc tuyến đầu (chẳng hạn như IDS/IPS hoạt động ở tốc độ đường truyền - wire-speed) vẫn mang bản chất tất định và thiếu linh hoạt. Mặc dù hiệu quả trong việc cắt giảm lưu lượng, chúng thiếu khả năng suy luận về các dấu hiệu cảnh báo mơ hồ hoặc cung cấp những diễn giải ngữ cảnh chuyên sâu. Hệ quả là, tốc độ dòng chảy của dữ liệu tạo ra một ``Nghịch lý Dữ liệu Lớn'' (Big Data Paradox)~\cite{bigdata2019}, đẩy gánh nặng nhận thức và phán đoán ngữ cảnh về phía các chuyên gia phân tích (analysts). Sự mất cân bằng cấu trúc này dẫn đến một cuộc khủng hoảng \emph{quá tải cảnh báo} (alert fatigue) nghiêm trọng, nơi các chuyên gia bị nhấn chìm bởi vô số thông báo có độ tin cậy thấp và phần lớn là báo động giả~\cite{alertfatigue2022}, làm gia tăng đáng kể nguy cơ bỏ lọt các cuộc xâm nhập thực sự.

Vấn đề nút thắt cổ chai này càng trở nên trầm trọng bởi hai lỗ hổng căn bản. Thứ nhất, các hệ thống lọc dựa trên chữ ký truyền thống~\cite{mlids2016} bộc lộ sự ``mù lòa'' về mặt cấu trúc trước các véc-tơ tấn công mới (như zero-day hoặc mã độc đa hình)---những mối đe dọa thường xuyên qua mặt các tập luật tĩnh chỉ bằng những sửa đổi nhỏ trên dữ liệu tải (payload). Điều này phơi bày nhu cầu cấp thiết về một cơ chế phát hiện dị thường thống kê không phụ thuộc vào chữ ký. Thứ hai, việc lấy mẫu log nhằm tối ưu hóa chi phí vô tình làm đứt gãy các liên kết thời gian cần thiết để tái dựng các chuỗi tấn công có tính chất chậm-và-lặng (low-and-slow) của các Chiến dịch Đe dọa Liên tục Bậc cao (APT)~\cite{apt_stealth2021}. Các nền tảng Điều phối, Tự động hóa và Phản hồi An ninh (SOAR) thông thường cũng không thể khỏa lấp khoảng trống suy luận này, do chúng hoạt động dựa trên các kịch bản (playbooks) cứng nhắc, không có khả năng suy luận thích ứng hay giải thích các cảnh báo nằm ngoài kịch bản định trước~\cite{soar2025}.

Các Mô hình Ngôn ngữ Lớn (LLM) hứa hẹn cung cấp chính xác năng lực suy luận ngữ cảnh chuyên sâu còn thiếu sót này; tuy nhiên, việc ứng dụng chúng trực tiếp vào môi trường SOC thực tế vấp phải một nghịch lý tích hợp đa chiều. Độ trễ suy luận của chúng là một trở ngại lớn đối với các luồng dữ liệu tốc độ cao, nguyên nhân cốt lõi từ chi phí tính toán bậc hai của cơ chế tự chú ý (self-attention)~\cite{vaswani2017}. Thêm vào đó, giao diện tương tác bằng ngôn ngữ tự nhiên khiến chúng đặc biệt dễ bị tổn thương trước các cuộc tấn công tiêm nhiễm câu lệnh qua nền tảng log (log-substrate prompt injection). Trong các kịch bản này, dữ liệu tải của kẻ tấn công được nhúng ngụy trang trong các trường thông tin (ví dụ: HTTP User-Agent) có thể đánh chiếm quá trình suy luận của mô hình nếu không có cơ chế cách ly nghiêm ngặt~\cite{promptinjection2023}. Cuối cùng, các quy định khắt khe về chủ quyền dữ liệu và mô hình Zero-Trust~\cite{nist80207} thường nghiêm cấm việc chuyển xuất dữ liệu telemetry nhạy cảm lên các dịch vụ đám mây bên ngoài, bắt buộc hệ thống suy luận phải được vận hành cục bộ (on-premises). Để giải quyết những thách thức này, nghiên cứu đề xuất SENTINEL: một kiến trúc nhận thức hai tầng. Kiến trúc này bao gồm một bộ lọc thống kê phi trạng thái (stateless) Tầng 1 hoạt động ở tốc độ đường truyền, được tăng cường bởi một Cổng Máy học (ML gateway) đảm nhiệm tự động giải quyết phần lớn các cảnh báo có độ tin cậy cao. Cấu phần này được tích hợp liền mạch với Tầng 2 AI Tác tử (Agentic AI) có trạng thái (stateful), được lưu trữ cục bộ, chuyên trách thực hiện các suy luận bậc cao trong môi trường cách ly nhận thức an toàn. Như một hệ quả tất yếu của cơ chế điều phối có trạng thái, thiết kế này tự động bảo tồn các chuỗi sự kiện theo thời gian và phát hiện các ngoại lai không có chữ ký, thiết lập một nền tảng vững chắc cho thế hệ hệ thống phòng thủ tự động tiếp theo.

\section{Mục tiêu Nghiên cứu}
Mục tiêu tổng quát của luận văn là thiết kế, triển khai, và đánh giá định lượng kiến trúc SENTINEL---một hệ thống nhận thức hai tầng kết hợp khả năng lọc thống kê độ trễ thấp với suy luận tác tử an toàn, nhận thức ngữ cảnh. Để hiện thực hóa điều này, nghiên cứu theo đuổi ba mục tiêu cụ thể:

\textbf{Mục tiêu 1 (Tầng 1 --- Hấp thụ lưu lượng):} Thiết kế một \emph{đường ống lọc hai chặng} nhằm giảm thiểu tình trạng quá tải cảnh báo ngay tại lớp thu nhận dữ liệu. Điều này bao gồm việc triển khai một bộ lọc tất định hoạt động ở tốc độ đường truyền, sử dụng thuật toán Welford trực tuyến để tính toán phương sai lăn và điểm Z-score với độ phức tạp $O(1)$, nối tiếp với một Cổng Máy học (LightGBM). Mục tiêu không nằm ở việc tối đa hóa độ chính xác phân loại đơn lẻ, mà là tối đa hóa khả năng giải quyết tự chủ các cảnh báo rõ ràng, qua đó cắt giảm đáng kể tần suất gọi đến cỗ máy suy luận LLM đòi hỏi tài nguyên tính toán cao.

\textbf{Mục tiêu 2 (Tầng 2 --- Đóng gói bảo mật cho LLM):} Xây dựng các cơ chế đóng gói bảo mật mang tính cấu trúc cho việc tích hợp mô hình LLM lượng tử hóa cục bộ. Yêu cầu này đòi hỏi thiết kế các cơ chế phòng thủ tất định dựa trên kiến trúc---cụ thể là Đóng gói Dữ liệu Phân định sử dụng Nonce mật mã (ví dụ: \texttt{<<<DATA\_HASH...>>>}), rào chắn đồng thuận hai tầng, cơ chế làm sạch đầu ra, và chuỗi kiểm toán sử dụng HMAC. Mục tiêu là vô hiệu hóa các véc-tơ đối kháng (như log-substrate prompt injection) bằng các bảo đảm về mặt kiến trúc thay vì phụ thuộc vào các bộ phân loại ngữ nghĩa vốn mong manh và dễ bị vượt qua.

\textbf{Mục tiêu 3 (Tầng 2 --- Điều phối và Nhận thức):} Xây dựng một khung điều phối nhận thức có trạng thái sử dụng LangGraph, tích hợp một Bộ nhớ Đe dọa (Threat Memory) bền vững với hệ thống sinh văn bản tăng cường truy xuất (RAG) kép. Bằng việc kết hợp tìm kiếm ngữ nghĩa dày đặc với đối sánh từ khóa thưa (MITRE ATT\&CK và NIST SP 800-61r2) thông qua thuật toán Hợp nhất Hạng Tương hỗ (Reciprocal Rank Fusion), mục tiêu này hướng tới việc tạo cơ sở vững chắc cho các quy kết pháp y và biện minh của tác tử, đồng thời cho phép khả năng tương quan nối kết các chiến dịch tấn công kéo dài nhiều ngày xuất phát từ trạng thái tích lũy.

\section{Câu hỏi Nghiên cứu}
\label{sec:rq}
Quá trình kiểm chứng thực nghiệm đối với kiến trúc được định hướng bởi ba câu hỏi nghiên cứu trọng tâm, mỗi câu đối ứng trực tiếp với các mục tiêu đã nêu:

\textbf{RQ1 (Kinh tế học của tầng lọc):} Một đường ống lai---bao gồm bộ lọc tất định $O(1)$ dựa trên đường cơ sở thống kê trực tuyến, kết hợp với một cổng ML học máy---có thể hấp thụ hiệu quả bao nhiêu phần trăm lưu lượng telemetry của SOC trước khi cần gọi đến tiến trình suy luận LLM cục bộ kéo dài nhiều giây, và sự đánh đổi định lượng về độ chính xác phát hiện là như thế nào?

\textbf{RQ2 (Ranh giới phòng thủ kiến trúc):} Trong khuôn khổ kiến trúc, việc bố trí cơ chế phòng thủ chống tiêm nhiễm nền-log (log-substrate injection) ở đâu mang lại hiệu quả cao nhất? Hơn nữa, năng lực của lớp lọc mẫu tất định so sánh như thế nào với cơ chế Đóng gói Dữ liệu Phân định (sử dụng nonces) và rào chắn đồng thuận hai tầng trong việc vô hiệu hóa các véc-tơ tấn công hạ cấp (downgrade attacks)?

\textbf{RQ3 (Hiệu năng của hệ thống có trạng thái):} Những năng lực đặc thù nào được kích hoạt khi chuyển đổi từ Tự động hóa Runbook tĩnh sang một khung tác tử có trạng thái được trang bị Bộ nhớ Đe dọa bền vững và cơ chế truy xuất tri thức kép? Hiệu quả được đo lường thông qua (a) khả năng nhận diện và tương quan các chiến dịch đa ngày mà không cần định nghĩa luật thời gian trước, và (b) độ chính xác trong việc quy kết kỹ thuật MITRE ATT\&CK cùng độ tin cậy của các biện minh sinh ra.

\begin{quote}\small
\textbf{Ghi chú về công cụ đo lường (Note on Instrumentation).} RQ3 chủ ý không sử dụng độ chính xác phân loại nhị phân làm thước đo cho đóng góp của hệ thống truy xuất tri thức. Ngữ cảnh truy xuất tác động căn bản tới khả năng \emph{quy kết} và \emph{biện minh} cho một phán quyết, thay vì xác định nhị phân tính chất độc hại. Việc áp dụng độ chính xác nhị phân trong bối cảnh này sẽ cấu thành một sai lầm về mặt phương pháp luận. Tương tự, luận văn này không khẳng định tầng LLM hoạt động như một bộ phân loại sơ cấp; các bằng chứng thực nghiệm định vị nó như một lưới an toàn thiên về độ phủ (coverage-biased) dành cho suy luận phức hợp, một giới hạn được công nhận rõ ràng ở Chương~\ref{ch:conclusion}.
\end{quote}

\section{Phạm vi và Đối tượng Nghiên cứu}
\textbf{Đối tượng Nghiên cứu:} Đối tượng chính của nghiên cứu là bản thân kiến trúc nhận thức hai tầng SENTINEL. Bao gồm các cơ chế vận hành chính xác yếu tố nào mà một bộ lọc thống kê Tầng 1 phi trạng thái ở tốc độ đường truyền và Cổng ML được điều phối cùng khung tác tử Tầng 2 có trạng thái, lưu trữ cục bộ, nhằm tự động hóa quy trình phát hiện và phản hồi của SOC. Các yếu tố này bao gồm: telemetry luồng dữ liệu tốc độ cao ở OSI Lớp 3--4 (NetFlow five-tuple và các đặc trưng CICFlowMeter, sử dụng siêu dữ liệu PCAP thay vì Phân tích Gói tin Sâu - Deep Packet Inspection); điều phối tác tử có trạng thái thông qua LangGraph; thực thi suy luận biệt lập sử dụng LLM lượng tử hóa cục bộ; và các biện pháp giảm thiểu véc-tơ tấn công như log-substrate injection, trốn tránh bộ phân cách (delimiter smuggling), và đầu độc cơ sở tri thức.

\textbf{Phạm vi Nghiên cứu:} Đánh giá thực nghiệm bị giới hạn nghiêm ngặt trong tập dữ liệu CSE-CIC-IDS2018~\cite{ids2018} cho phân loại tấn công diện rộng, và tập dữ liệu DAPT2020~\cite{dapt2020} cho tương quan APT đa giai đoạn. Các mối đe dọa Zero-day được mô phỏng bằng cách tiêm các độ lệch thống kê Welford cực đoan vào các luồng dữ liệu an toàn đã được xác minh. Cỗ máy suy luận được chuẩn hóa trên mô hình \textit{Gemma-2-9B-IT} chạy cục bộ, được lượng tử hóa theo định dạng Q6\_K qua \texttt{llama.cpp}, thực thi trên một GPU tiêu dùng (consumer-grade) duy nhất nhằm đảm bảo tuân thủ tuyệt đối về chủ quyền dữ liệu. Thêm vào đó, năng lực phản hồi tự chủ được giới hạn ở việc sinh các chính sách tường lửa, báo cáo pháp y, và quy kết cấu trúc; việc tự động ngắt kết nối vật lý bị loại trừ rõ ràng khỏi phạm vi hiện tại.

\section{Phương pháp Nghiên cứu}
Luận văn này áp dụng phương pháp Nghiên cứu Khoa học Thiết kế (Design Science Research - DSR), được chống đỡ bởi lập luận suy diễn và định lượng. Nó kiến tạo tri thức mới bằng cách xây dựng một công cụ thực tế (artifact) hoạt động (SENTINEL) và đánh giá nghiêm ngặt hiệu năng của nó trước các thách thức vận hành thực tiễn, đảm bảo mọi tuyên bố lý thuyết đều được chứng minh bằng thực nghiệm. Khung suy luận hoạt động trên ba bình diện: thống kê (phân định đường cơ sở của lưu lượng bình thường so với sự bất thường), logic (chính thức hóa đồ thị trạng thái ra quyết định của tác tử), và đối kháng (kiểm tra ứng suất khả năng phục hồi của hệ thống trước sự thao túng).

Về mặt thực tiễn, kiến trúc được triển khai dưới dạng một nguyên mẫu hai tầng hoàn chỉnh. Nguyên mẫu này trải qua quá trình đánh giá từ đầu đến cuối (end-to-end) bằng cách phát lại các tập dữ liệu lưu trữ và lưu lượng trực tiếp mô phỏng, trong đó có chủ đích đưa vào các tải dữ liệu đối kháng để đánh giá tính mạnh mẽ. Hệ thống được đánh giá qua một khung năm chiều toàn diện: Độ chính xác, Hiệu suất, Bảo mật, Tính Giải thích, và Tính Toàn vẹn. Các nghiên cứu bóc tách (ablation studies) có kiểm soát được sử dụng để cô lập hệ thống đóng góp của từng cấu phần riêng lẻ, trong khi các kiểm định thống kê phi tham số (ví dụ: McNemar và Mann--Whitney U) được áp dụng nhằm xác nhận ý nghĩa thống kê của các biến thể kiến trúc. Quá trình đánh giá tổng hợp các số liệu định lượng (như độ trễ, điểm F1) với các kết quả định tính (lời biện minh pháp y của mô hình). Để ngăn ngừa xu hướng tự tôn (self-aggrandisement bias), khả năng suy luận định tính được chấm điểm khách quan thông qua một khung tự động lấy cảm hứng từ RAGAS~\cite{ragas2023}, do một trọng tài LLM chéo họ (cross-family) độc lập điều hành, đảm bảo một đánh giá so sánh nghiêm ngặt.

\section{Đóng góp của Luận văn}
\label{sec:contributions}
Các đóng góp của nghiên cứu này bao gồm ba khía cạnh, mỗi khía cạnh được phân định rạch ròi bởi các bằng chứng thực nghiệm đi kèm.

\textbf{Về Kiến trúc,} luận văn xác thực một đường ống lai hai tầng giải quyết thành công các vấn đề thắt nút cổ chai về độ trễ liên quan đến việc triển khai các LLM cục bộ (vốn đòi hỏi thời gian suy luận kéo dài nhiều giây) trong môi trường SOC vận hành thực tế. Đóng góp cốt lõi là \emph{tích hợp có tính bố cục (compositional integration)}: một bộ lọc thống kê $O(1)$ phi trạng thái kết hợp với một Cổng LightGBM có khả năng tự chủ cắt bỏ $83.8\%$ các trường hợp cảnh báo rõ ràng ở tốc độ đường truyền. Hơn thế nữa, một Semantic Cache bộ nhớ trong (Tầng 1.75) giúp giải quyết các cảnh báo trùng lặp chỉ trong khoảng $1$\,ms. Cơ chế này đảm bảo mạng lưới được bảo vệ một cách tất định, dành riêng tầng nhận thức (vốn đắt đỏ về tính toán) cho một thiểu số nhỏ các sự kiện bất thường mơ hồ.

\textbf{Về Bảo mật AI,} nghiên cứu đóng góp một cơ chế phòng thủ tất định, mạnh mẽ về mặt cấu trúc chống lại các cuộc tấn công log-substrate prompt injection. Thông qua việc áp dụng Đóng gói Dữ liệu Phân định sử dụng cryptographic nonces (ví dụ: \texttt{<<<DATA\_HASH...>>>}), một rào chắn đồng thuận Tầng 1/Tầng 2, và chuỗi kiểm toán HMAC, hệ thống cung cấp các bảo đảm an ninh có giá trị thông qua kiến trúc cấu trúc thay vì dựa trên các mô hình huấn luyện mang tính xác suất. Điều này đảm bảo rằng các luật cố định của rào chắn đồng thuận không thể bị vượt qua thông qua sự thuyết phục, qua đó vô hiệu hóa hiệu quả các nỗ lực tấn công hạ cấp (downgrade) cả về cú pháp lẫn ngữ nghĩa.

\textbf{Về Điều phối Tác tử,} kiến trúc thay thế Tự động hóa Runbook cứng nhắc bằng một đồ thị trạng thái có trạng thái, không chu trình (acyclic), và không phụ thuộc công cụ (tool-free), được điều phối trên một mô hình lượng tử hóa cục bộ duy nhất. Điều này tạo ra một tầng phản hồi có khả năng thích ứng cao, nhưng vẫn đảm bảo tính lặp lại nghiêm ngặt và khả năng kiểm toán pháp y toàn vẹn, hoàn toàn tuân thủ các ràng buộc khắt khe về chủ quyền dữ liệu. Các năng lực nền tảng bổ sung được kích hoạt bởi kiến trúc này---cụ thể là leo thang sự cố dị thường không cần chữ ký thông qua thống kê đường cơ sở, và tương quan sự kiện APT đa ngày từ trạng thái tích lũy---được trình bày như những bằng chứng khái niệm (proof of concept) trong giới hạn, đặt nền tảng cho việc thương mại hóa và mở rộng sản xuất trong tương lai.

\section{Cấu trúc của Luận văn}
Chương~\ref{ch:theoretical_background} thiết lập nền tảng lý thuyết, bao gồm lượng tử hóa LLM, mô hình đồ thị trạng thái, RAG lai, thống kê trực tuyến, và tổng quan các tài liệu về an ninh tác tử liên quan. Chương~\ref{ch:system_design_implementation} đi sâu vào thiết kế hệ thống kỹ thuật: bộ lọc hai tầng, Cổng ML, Máy Trạng thái Hữu hạn (FSM) tác tử có trạng thái, truy xuất lai, rào chắn cấu trúc, bảng điều khiển (dashboard), và các cấu hình triển khai. Chương~\ref{ch:experiments_evaluation} trình bày các thí nghiệm thực chứng, các nghiên cứu bóc tách có kiểm soát, và kiểm chứng thống kê phi tham số. Cuối cùng, Chương~\ref{ch:conclusion} tổng hợp các phát hiện, nhìn nhận các hạn chế hiện tại, và phác thảo các định hướng nghiên cứu trong tương lai.
